% GENERATED FILE. Edit canonical lessons, then run npm run generate:books.
% canonical-source-hash: fnv1a64:99ebb0097734b6d4
% canonical-lessons: KA-C439-visa, KA-C439-kastams, KA-C439-sunka, KA-C439-vinimaya-madu, KA-C439-etiem

\chapter{Crossing a Border}
\label{ch:ka-crossing-a-border}
\hlchaptermodality{439}

\begin{chapteropening}
\textbf{By the end of this chapter:} \emph{I can say a visa, customs, a customs duty, to exchange money and a cash machine.}

Complete the last lesson of chapter 439: I can say a visa, customs, a customs duty, to exchange money and a cash machine.
\end{chapteropening}

\section[\texorpdfstring{\kn{ವೀಸಾ} (vīsā)}{ವೀಸಾ (vīsā)}]{\kn{ವೀಸಾ} (vīsā) — a visa}

\label{lesson:KA-C439-visa}



\begin{quote}
Before the new one: say the Kannada for air travel, then the Kannada for a helicopter.
\end{quote}

\subsection*{You'll want to know: \kn{ವೀಸಾ}}
\textbf{\kn{ವೀಸಾ}} — \emph{vīsā} — \textquotedblleft{}a visa\textquotedblright{}.

\begin{quote}
\kn{ವಿದೇಶಕ್ಕೆ} \kn{ಹೋಗಲು} \kn{ಪಾಸ್ಪೋರ್ಟ್} \kn{ಮತ್ತು} \kn{ವೀಸಾ} \kn{ಬೇಕು}.
\end{quote}

\begin{quote}
\emph{vidēśakke hōgalu pāspōrṭ mattu vīsā bēku.}
\end{quote}

\begin{quote}
To go abroad you need a passport and a visa.
\end{quote}

\kn{ವೀಸಾ} \kn{ಸಿಕ್ಕಿತು}, the visa came through.

\begin{cousinweb}
\textbf{\kn{ವೀಸಾ}} (\emph{vīsā}) is a loan from English \emph{visa}.
\end{cousinweb}

\begin{grammarlens}[title={What you've built}]
One more word for getting around.
\end{grammarlens}

\subsection*{Your turn}
\begin{itemize}
\raggedright
  \item \emph{Say it:} \emph{vīsā}
  \item \emph{Say it:} \emph{vīsā}, once more
  \item \emph{Say it:} \emph{vidēśakke hōgalu pāspōrṭ mattu vīsā bēku.} — To go abroad you need a passport and a visa.
  \item \emph{Recall:} say the Kannada for air travel, then the Kannada for a helicopter, then say \emph{vīsā} again
\end{itemize}

\begin{tcolorbox}[breakable,colback=teal!4,colframe=teal!35!black,title={Before you move on}]
What does \kn{ವೀಸಾ} mean? (\textquotedblleft{}A visa\textquotedblright{}.) Say it once more.
\end{tcolorbox}
\section[\texorpdfstring{\kn{ಕಸ್ಟಮ್ಸ್} (kasṭams)}{ಕಸ್ಟಮ್ಸ್ (kasṭams)}]{\kn{ಕಸ್ಟಮ್ಸ್} (kasṭams) — customs}

\label{lesson:KA-C439-kastams}



\begin{quote}
Before the new one: say the Kannada for a helicopter, then the Kannada for a visa.
\end{quote}

\subsection*{You'll want to know: \kn{ಕಸ್ಟಮ್ಸ್}}
\textbf{\kn{ಕಸ್ಟಮ್ಸ್}} — \emph{kasṭams} — \textquotedblleft{}customs\textquotedblright{}.

\begin{quote}
\kn{ಕಸ್ಟಮ್ಸ್} \kn{ಅಧಿಕಾರಿಗಳು} \kn{ನಮ್ಮ} \kn{ಚೀಲಗಳನ್ನು} \kn{ತೆರೆದರು}.
\end{quote}

\begin{quote}
\emph{kasṭams adhikārigaḷu namma cīlagaḷannu teredaru.}
\end{quote}

\begin{quote}
The customs officers opened our bags.
\end{quote}

The office at a border or an airport that checks goods.

\begin{cousinweb}
\textbf{\kn{ಕಸ್ಟಮ್ಸ್}} (\emph{kasṭams}) is a loan from English \emph{customs}, the duties charged on goods and the office that collects them.
\end{cousinweb}

\begin{grammarlens}[title={What you've built}]
One more word for getting around.
\end{grammarlens}

\subsection*{Your turn}
\begin{itemize}
\raggedright
  \item \emph{Say it:} \emph{kasṭams}
  \item \emph{Say it:} \emph{kasṭams}, once more
  \item \emph{Say it:} \emph{kasṭams adhikārigaḷu namma cīlagaḷannu teredaru.} — The customs officers opened our bags.
  \item \emph{Recall:} say the Kannada for a helicopter, then the Kannada for a visa, then say \emph{kasṭams} again
\end{itemize}

\begin{tcolorbox}[breakable,colback=teal!4,colframe=teal!35!black,title={Before you move on}]
What does \kn{ಕಸ್ಟಮ್ಸ್} mean? (\textquotedblleft{}Customs\textquotedblright{}.) Say it once more.
\end{tcolorbox}
\section[\texorpdfstring{\kn{ಸುಂಕ} (suṅka)}{ಸುಂಕ (suṅka)}]{\kn{ಸುಂಕ} (suṅka) — a customs duty; a toll}

\label{lesson:KA-C439-sunka}



\begin{quote}
Before the new one: say the Kannada for a visa, then the Kannada for customs.
\end{quote}

\subsection*{You'll want to know: \kn{ಸುಂಕ}}
\textbf{\kn{ಸುಂಕ}} — \emph{suṅka} — \textquotedblleft{}a customs duty; a toll\textquotedblright{}.

\begin{quote}
\kn{ಚಿನ್ನದ} \kn{ಆಭರಣಗಳಿಗೆ} \kn{ಸುಂಕ} \kn{ಕಟ್ಟಬೇಕು}.
\end{quote}

\begin{quote}
\emph{cinnada ābharaṇagaḷige suṅka kaṭṭabēku.}
\end{quote}

\begin{quote}
Duty has to be paid on gold jewellery.
\end{quote}

\kn{ಸುಂಕ} \kn{ಕಟ್ಟು}, to pay duty. \kn{ಶುಲ್ಕ}, a fee, is the learned form of the same Sanskrit word.

\begin{cousinweb}
\textbf{\kn{ಸುಂಕ}} (\emph{suṅka}) comes from Sanskrit \emph{śulka}, \textquotedblleft{}price, toll\textquotedblright{}, through a Prakrit form \emph{suṅka} (CDIAL 12543). \kn{ಶುಲ್ಕ}, a fee, is the same word borrowed again in its learned shape.
\end{cousinweb}

\begin{grammarlens}[title={What you've built}]
One more word for getting around.
\end{grammarlens}

\subsection*{Your turn}
\begin{itemize}
\raggedright
  \item \emph{Say it:} \emph{suṅka}
  \item \emph{Say it:} \emph{suṅka}, once more
  \item \emph{Say it:} \emph{cinnada ābharaṇagaḷige suṅka kaṭṭabēku.} — Duty has to be paid on gold jewellery.
  \item \emph{Recall:} say the Kannada for a visa, then the Kannada for customs, then say \emph{suṅka} again
\end{itemize}

\begin{tcolorbox}[breakable,colback=teal!4,colframe=teal!35!black,title={Before you move on}]
What does \kn{ಸುಂಕ} mean? (\textquotedblleft{}A customs duty; a toll\textquotedblright{}.) Say it once more.
\end{tcolorbox}
\section[\texorpdfstring{\kn{ವಿನಿಮಯ} \kn{ಮಾಡು} (vinimaya māḍu)}{ವಿನಿಮಯ ಮಾಡು (vinimaya māḍu)}]{\kn{ವಿನಿಮಯ} \kn{ಮಾಡು} (vinimaya māḍu) — to exchange money; to swap}

\label{lesson:KA-C439-vinimaya-madu}



\begin{quote}
Before the new one: say the Kannada for customs, then the Kannada for a customs duty.
\end{quote}

\subsection*{You'll want to know: \kn{ವಿನಿಮಯ} \kn{ಮಾಡು}}
\textbf{\kn{ವಿನಿಮಯ} \kn{ಮಾಡು}} — \emph{vinimaya māḍu} — \textquotedblleft{}to exchange money; to swap\textquotedblright{}.

\begin{quote}
\kn{ವಿಮಾನ} \kn{ನಿಲ್ದಾಣದಲ್ಲಿ} \kn{ನಾನು} \kn{ಡಾಲರನ್ನು} \kn{ರೂಪಾಯಿಗೆ} \kn{ವಿನಿಮಯ} \kn{ಮಾಡಿದೆ}.
\end{quote}

\begin{quote}
\emph{vimāna nildāṇadalli nānu ḍālarannu rūpāyige vinimaya māḍide.}
\end{quote}

\begin{quote}
At the airport I changed dollars into rupees.
\end{quote}

Past (I …): \kn{ವಿನಿಮಯ} \kn{ಮಾಡಿದೆ}, \emph{vinimaya māḍide}.

\kn{ವಿನಿಮಯ} \kn{ದರ}, an exchange rate; what you get takes -\kn{ಗೆ}.

\begin{cousinweb}
\textbf{\kn{ವಿನಿಮಯ} \kn{ಮಾಡು}} (\emph{vinimaya māḍu}) joins Sanskrit \emph{vinimaya}, \textquotedblleft{}exchange, barter\textquotedblright{} (Monier-Williams), and the native verb \kn{ಮಾಡು}, to do.
\end{cousinweb}

\begin{grammarlens}[title={What you've built}]
One more word for getting around.
\end{grammarlens}

\subsection*{Your turn}
\begin{itemize}
\raggedright
  \item \emph{Say it:} \emph{vinimaya māḍu}
  \item \emph{Say it:} \emph{vinimaya māḍu}, once more
  \item \emph{Say it:} \emph{vimāna nildāṇadalli nānu ḍālarannu rūpāyige vinimaya māḍide.} — At the airport I changed dollars into rupees.
  \item \emph{Recall:} say the Kannada for customs, then the Kannada for a customs duty, then say \emph{vinimaya māḍu} again
\end{itemize}

\begin{tcolorbox}[breakable,colback=teal!4,colframe=teal!35!black,title={Before you move on}]
What does \kn{ವಿನಿಮಯ} \kn{ಮಾಡು} mean? (\textquotedblleft{}To exchange money; to swap\textquotedblright{}.) Say it once more.
\end{tcolorbox}
\section[\texorpdfstring{\kn{ಎಟಿಎಂ} (eṭieṃ)}{ಎಟಿಎಂ (eṭieṃ)}]{\kn{ಎಟಿಎಂ} (eṭieṃ) — a cash machine; an ATM}

\label{lesson:KA-C439-etiem}



\begin{quote}
Before the new one: say the Kannada for a customs duty, then the Kannada for to exchange money.
\end{quote}

\subsection*{You'll want to know: \kn{ಎಟಿಎಂ}}
\textbf{\kn{ಎಟಿಎಂ}} — \emph{eṭieṃ} — \textquotedblleft{}a cash machine; an ATM\textquotedblright{}.

\begin{quote}
\kn{ನಾನು} \kn{ಎಟಿಎಂನಿಂದ} \kn{ಸ್ವಲ್ಪ} \kn{ದುಡ್ಡು} \kn{ತೆಗೆದುಕೊಂಡೆ}.
\end{quote}

\begin{quote}
\emph{nānu eṭieṃninda svalpa duḍḍu tegedukoṇḍe.}
\end{quote}

\begin{quote}
I took some money out of the cash machine.
\end{quote}

Said letter by letter, as in English.

\begin{cousinweb}
\textbf{\kn{ಎಟಿಎಂ}} (\emph{eṭieṃ}) is a loan from English \emph{ATM}, the letters read out one by one.
\end{cousinweb}

\begin{grammarlens}[title={What you've built}]
One more word for getting around.
\end{grammarlens}

\subsection*{Your turn}
\begin{itemize}
\raggedright
  \item \emph{Say it:} \emph{eṭieṃ}
  \item \emph{Say it:} \emph{eṭieṃ}, once more
  \item \emph{Say it:} \emph{nānu eṭieṃninda svalpa duḍḍu tegedukoṇḍe.} — I took some money out of the cash machine.
  \item \emph{Recall:} say the Kannada for a customs duty, then the Kannada for to exchange money, then say \emph{eṭieṃ} again
\end{itemize}

\begin{tcolorbox}[breakable,colback=teal!4,colframe=teal!35!black,title={Before you move on}]
What does \kn{ಎಟಿಎಂ} mean? (\textquotedblleft{}A cash machine; an ATM\textquotedblright{}.) Say it once more.
\end{tcolorbox}
